
\documentclass{article}
\usepackage{colortbl}
\usepackage{makecell}
\usepackage{multirow}
\usepackage{supertabular}

\begin{document}

\newcounter{utterance}

\centering \large Interaction Transcript for game `function\_detective', experiment `two\_numbers\_active\_pair\_checks', episode 278 with Qwen3.6{-}35B{-}A3B{-}FP8{-}without{-}reasoning.
\vspace{24pt}

{ \footnotesize  \setcounter{utterance}{1}
\setlength{\tabcolsep}{0pt}
\begin{supertabular}{c@{$\;$}|p{.15\linewidth}@{}p{.15\linewidth}p{.15\linewidth}p{.15\linewidth}p{.15\linewidth}p{.15\linewidth}}
    \# & \multicolumn{2}{c}{Player} && \multicolumn{2}{c}{Game Master} \\
    \hline

    \theutterance \stepcounter{utterance}  
    & & & \multicolumn{4}{p{0.6\linewidth}}{
        \cellcolor[rgb]{0.9,0.9,0.9}{
            \makecell[{{p{\linewidth}}}]{
                \texttt{\tiny{[P1$\langle$GM]}}
                \texttt{I have a hidden Python function for you to figure out.} \\
\\ 
\texttt{Your job is to infer what it does and then submit a correct Python implementation.} \\
\\ 
\texttt{Function signature:} \\
\texttt{(a: int, b: int) {-}> int} \\
\\ 
\texttt{You have 15 turns total.} \\
\\ 
\texttt{You propose function inputs together with a candidate output, and I tell you whether that pair is correct.} \\
\\ 
\texttt{You must make your best SOLVE attempt no later than the second{-}to{-}last turn.} \\
\texttt{Do not wait until the very end.} \\
\\ 
\texttt{You can use these actions:} \\
\texttt{{-} JSON format described at the end of the prompt} \\
\texttt{{-} SOLVE: ```python ... ```} \\
\\ 
\texttt{SOLVE: ```python} \\
\texttt{def solution(<value1>, <value2>):} \\
\texttt{    ...} \\
\texttt{```} \\
\\ 
\texttt{Do not add import statements in your guess.} \\
\texttt{Anything you might need is already available, so just write the function body and return the answer.} \\
\\ 
\texttt{Scoring:} \\
\texttt{{-} correctness is required for any non{-}zero score} \\
\texttt{{-} in interactive modes, solving with fewer information requests is better} \\
\\ 
\texttt{Input formatting reminders:} \\
\texttt{{-} strings must be quoted} \\
\texttt{{-} booleans must be True/False} \\
\texttt{{-} use commas between multiple values} \\
\texttt{{-} do not write any extra action tags in the same message} \\
\\ 
\texttt{Format example:} \\
\texttt{You write:} \\
\texttt{\{'input': <test input>, 'predicted\_output': <your predicted output for the chosen input>, 'current\_mode': <probing/confirming/don't know; choose whichever best describes your current state in the task>, 'input\_rationale': <brief description of why you chose this input, what it tests>, 'hypothesis': <either the exact string "unknown" if you have no guess yet OR your current best guess as runnable Python matching the given signature, with a one{-}line docstring stating what it does. Never write placeholder code, if you cannot commit to a guess, use "unknown". Match the given signature.} \\
\texttt{e.g.} \\
\texttt{def f(x):} \\
\texttt{	"""Returns the absolute value of x."""} \\
\texttt{	return x if x >= 0 else {-}x} \\
\\ 
\texttt{Use "unknown" if you have no guess yet.>, 'confidence\_score': <score between 1{-}100. 100 means highest confidence in your hypothesis>\}} \\
\\ 
\texttt{I answer:} \\
\texttt{OUTPUT: True} \\
\\ 
\texttt{Output only one valid json per turn.} \\
            }
        }
    }
     \\ \\

    \theutterance \stepcounter{utterance}  
    & \multicolumn{4}{p{0.6\linewidth}}{
        \cellcolor[rgb]{0.561,0.737,0.561}{
            \makecell[{{p{\linewidth}}}]{
                \texttt{\tiny{[P1$\rangle$GM]}}
                \texttt{\{'input': {[}1, 2{]}, 'predicted\_output': 3, 'current\_mode': 'probing', 'input\_rationale': 'Starting with simple small integers to understand the basic behavior of the function. Testing if it might be addition.', 'hypothesis': 'unknown', 'confidence\_score': 10\}} \\
            }
        }
    }
    & & \\ \\

    \theutterance \stepcounter{utterance}  
    & & & \multicolumn{4}{p{0.6\linewidth}}{
        \cellcolor[rgb]{0.9,0.9,0.9}{
            \makecell[{{p{\linewidth}}}]{
                \texttt{\tiny{[P1$\langle$GM]}}
                \texttt{Game Violation: ParseError: Response is not valid JSON.} \\
\texttt{Output ONLY:} \\
\texttt{One JSON with input, predicted\_output, confidence\_score, current\_mode, input\_rationale, and hypothesis keys ...} \\
\texttt{OR} \\
\texttt{SOLVE: ```python ... ```} \\
            }
        }
    }
     \\ \\

    \theutterance \stepcounter{utterance}  
    & \multicolumn{4}{p{0.6\linewidth}}{
        \cellcolor[rgb]{0.561,0.737,0.561}{
            \makecell[{{p{\linewidth}}}]{
                \texttt{\tiny{[P1$\rangle$GM]}}
                \texttt{\{} \\
\texttt{  "input": {[}1, 2{]},} \\
\texttt{  "predicted\_output": 3,} \\
\texttt{  "current\_mode": "probing",} \\
\texttt{  "input\_rationale": "Starting with simple small integers to understand the basic behavior of the function. Testing if it might be addition.",} \\
\texttt{  "hypothesis": "unknown",} \\
\texttt{  "confidence\_score": 10} \\
\texttt{\}} \\
            }
        }
    }
    & & \\ \\

    \theutterance \stepcounter{utterance}  
    & & & \multicolumn{4}{p{0.6\linewidth}}{
        \cellcolor[rgb]{0.9,0.9,0.9}{
            \makecell[{{p{\linewidth}}}]{
                \texttt{\tiny{[P1$\langle$GM]}}
                \texttt{OUTPUT: True} \\
            }
        }
    }
     \\ \\

    \theutterance \stepcounter{utterance}  
    & \multicolumn{4}{p{0.6\linewidth}}{
        \cellcolor[rgb]{0.561,0.737,0.561}{
            \makecell[{{p{\linewidth}}}]{
                \texttt{\tiny{[P1$\rangle$GM]}}
                \texttt{\{} \\
\texttt{  "input": {[}2, 3{]},} \\
\texttt{  "predicted\_output": 5,} \\
\texttt{  "current\_mode": "probing",} \\
\texttt{  "input\_rationale": "Testing another pair of small integers to confirm if the function is simply addition. If f(1,2)=3 and f(2,3)=5, it strongly suggests addition.",} \\
\texttt{  "hypothesis": "unknown",} \\
\texttt{  "confidence\_score": 20} \\
\texttt{\}} \\
            }
        }
    }
    & & \\ \\

    \theutterance \stepcounter{utterance}  
    & & & \multicolumn{4}{p{0.6\linewidth}}{
        \cellcolor[rgb]{0.9,0.9,0.9}{
            \makecell[{{p{\linewidth}}}]{
                \texttt{\tiny{[P1$\langle$GM]}}
                \texttt{OUTPUT: True} \\
            }
        }
    }
     \\ \\

    \theutterance \stepcounter{utterance}  
    & \multicolumn{4}{p{0.6\linewidth}}{
        \cellcolor[rgb]{0.561,0.737,0.561}{
            \makecell[{{p{\linewidth}}}]{
                \texttt{\tiny{[P1$\rangle$GM]}}
                \texttt{\{} \\
\texttt{  "input": {[}0, 0{]},} \\
\texttt{  "predicted\_output": 0,} \\
\texttt{  "current\_mode": "probing",} \\
\texttt{  "input\_rationale": "Testing edge case with zeros to ensure it behaves like addition (0+0=0) and not something else like multiplication (0*0=0) or bitwise OR (0|0=0). Since previous tests were consistent with addition, this helps rule out other operations if the result is non{-}zero, but here it confirms the additive identity.",} \\
\texttt{  "hypothesis": "unknown",} \\
\texttt{  "confidence\_score": 30} \\
\texttt{\}} \\
            }
        }
    }
    & & \\ \\

    \theutterance \stepcounter{utterance}  
    & & & \multicolumn{4}{p{0.6\linewidth}}{
        \cellcolor[rgb]{0.9,0.9,0.9}{
            \makecell[{{p{\linewidth}}}]{
                \texttt{\tiny{[P1$\langle$GM]}}
                \texttt{OUTPUT: True} \\
            }
        }
    }
     \\ \\

    \theutterance \stepcounter{utterance}  
    & \multicolumn{4}{p{0.6\linewidth}}{
        \cellcolor[rgb]{0.561,0.737,0.561}{
            \makecell[{{p{\linewidth}}}]{
                \texttt{\tiny{[P1$\rangle$GM]}}
                \texttt{\{} \\
\texttt{  "input": {[}1, 1{]},} \\
\texttt{  "predicted\_output": 2,} \\
\texttt{  "current\_mode": "probing",} \\
\texttt{  "input\_rationale": "Testing with identical inputs to further confirm addition behavior. If f(1,1)=2, it is consistent with addition. This also helps distinguish from operations like XOR (1\^{}1=0) or multiplication (1*1=1).",} \\
\texttt{  "hypothesis": "unknown",} \\
\texttt{  "confidence\_score": 40} \\
\texttt{\}} \\
            }
        }
    }
    & & \\ \\

    \theutterance \stepcounter{utterance}  
    & & & \multicolumn{4}{p{0.6\linewidth}}{
        \cellcolor[rgb]{0.9,0.9,0.9}{
            \makecell[{{p{\linewidth}}}]{
                \texttt{\tiny{[P1$\langle$GM]}}
                \texttt{OUTPUT: True} \\
            }
        }
    }
     \\ \\

    \theutterance \stepcounter{utterance}  
    & \multicolumn{4}{p{0.6\linewidth}}{
        \cellcolor[rgb]{0.561,0.737,0.561}{
            \makecell[{{p{\linewidth}}}]{
                \texttt{\tiny{[P1$\rangle$GM]}}
                \texttt{\{} \\
\texttt{  "input": {[}10, 20{]},} \\
\texttt{  "predicted\_output": 30,} \\
\texttt{  "current\_mode": "probing",} \\
\texttt{  "input\_rationale": "Testing with larger distinct integers to further confirm the addition hypothesis. If f(10,20)=30, it strongly supports that the function is addition.",} \\
\texttt{  "hypothesis": "unknown",} \\
\texttt{  "confidence\_score": 50} \\
\texttt{\}} \\
            }
        }
    }
    & & \\ \\

    \theutterance \stepcounter{utterance}  
    & & & \multicolumn{4}{p{0.6\linewidth}}{
        \cellcolor[rgb]{0.9,0.9,0.9}{
            \makecell[{{p{\linewidth}}}]{
                \texttt{\tiny{[P1$\langle$GM]}}
                \texttt{OUTPUT: True} \\
            }
        }
    }
     \\ \\

    \theutterance \stepcounter{utterance}  
    & \multicolumn{4}{p{0.6\linewidth}}{
        \cellcolor[rgb]{0.561,0.737,0.561}{
            \makecell[{{p{\linewidth}}}]{
                \texttt{\tiny{[P1$\rangle$GM]}}
                \texttt{\{} \\
\texttt{  "input": {[}{-}1, 1{]},} \\
\texttt{  "predicted\_output": 0,} \\
\texttt{  "current\_mode": "probing",} \\
\texttt{  "input\_rationale": "Testing with negative numbers to ensure the function handles them correctly as addition. If f({-}1, 1) = 0, it confirms addition behavior with negative inputs.",} \\
\texttt{  "hypothesis": "unknown",} \\
\texttt{  "confidence\_score": 60} \\
\texttt{\}} \\
            }
        }
    }
    & & \\ \\

    \theutterance \stepcounter{utterance}  
    & & & \multicolumn{4}{p{0.6\linewidth}}{
        \cellcolor[rgb]{0.9,0.9,0.9}{
            \makecell[{{p{\linewidth}}}]{
                \texttt{\tiny{[P1$\langle$GM]}}
                \texttt{OUTPUT: False} \\
            }
        }
    }
     \\ \\

    \theutterance \stepcounter{utterance}  
    & \multicolumn{4}{p{0.6\linewidth}}{
        \cellcolor[rgb]{0.561,0.737,0.561}{
            \makecell[{{p{\linewidth}}}]{
                \texttt{\tiny{[P1$\rangle$GM]}}
                \texttt{\{} \\
\texttt{  "input": {[}{-}1, 1{]},} \\
\texttt{  "predicted\_output": 0,} \\
\texttt{  "current\_mode": "probing",} \\
\texttt{  "input\_rationale": "The previous hypothesis of simple addition was incorrect because f({-}1, 1) != 0. This suggests the function might involve absolute values, bitwise operations, or some other logic that doesn't simply sum negative and positive numbers. Let's test with positive numbers again to see if the pattern holds for them, and try a case where the result would differ from addition if it were something like max(a,b) or abs(a{-}b).",} \\
\texttt{  "hypothesis": "unknown",} \\
\texttt{  "confidence\_score": 40} \\
\texttt{\}} \\
            }
        }
    }
    & & \\ \\

    \theutterance \stepcounter{utterance}  
    & & & \multicolumn{4}{p{0.6\linewidth}}{
        \cellcolor[rgb]{0.9,0.9,0.9}{
            \makecell[{{p{\linewidth}}}]{
                \texttt{\tiny{[P1$\langle$GM]}}
                \texttt{OUTPUT: False} \\
            }
        }
    }
     \\ \\

    \theutterance \stepcounter{utterance}  
    & \multicolumn{4}{p{0.6\linewidth}}{
        \cellcolor[rgb]{0.561,0.737,0.561}{
            \makecell[{{p{\linewidth}}}]{
                \texttt{\tiny{[P1$\rangle$GM]}}
                \texttt{\{} \\
\texttt{  "input": {[}1, 2{]},} \\
\texttt{  "predicted\_output": 3,} \\
\texttt{  "current\_mode": "probing",} \\
\texttt{  "input\_rationale": "Re{-}testing the initial case to ensure consistency. If f(1,2) is still 3, it suggests the function might be addition for positive numbers but behaves differently for negatives or zero. Or perhaps my previous test of {[}{-}1, 1{]} was misinterpreted. Let's try another positive pair to be sure.",} \\
\texttt{  "hypothesis": "unknown",} \\
\texttt{  "confidence\_score": 40} \\
\texttt{\}} \\
            }
        }
    }
    & & \\ \\

    \theutterance \stepcounter{utterance}  
    & & & \multicolumn{4}{p{0.6\linewidth}}{
        \cellcolor[rgb]{0.9,0.9,0.9}{
            \makecell[{{p{\linewidth}}}]{
                \texttt{\tiny{[P1$\langle$GM]}}
                \texttt{OUTPUT: True} \\
            }
        }
    }
     \\ \\

    \theutterance \stepcounter{utterance}  
    & \multicolumn{4}{p{0.6\linewidth}}{
        \cellcolor[rgb]{0.561,0.737,0.561}{
            \makecell[{{p{\linewidth}}}]{
                \texttt{\tiny{[P1$\rangle$GM]}}
                \texttt{\{} \\
\texttt{  "input": {[}1, {-}1{]},} \\
\texttt{  "predicted\_output": 0,} \\
\texttt{  "current\_mode": "probing",} \\
\texttt{  "input\_rationale": "Testing with a positive and a negative number to distinguish between addition (1 + ({-}1) = 0) and other operations like absolute difference (|1 {-} ({-}1)| = 2) or max (max(1, {-}1) = 1). This helps narrow down the function's behavior with mixed signs.",} \\
\texttt{  "hypothesis": "unknown",} \\
\texttt{  "confidence\_score": 40} \\
\texttt{\}} \\
            }
        }
    }
    & & \\ \\

    \theutterance \stepcounter{utterance}  
    & & & \multicolumn{4}{p{0.6\linewidth}}{
        \cellcolor[rgb]{0.9,0.9,0.9}{
            \makecell[{{p{\linewidth}}}]{
                \texttt{\tiny{[P1$\langle$GM]}}
                \texttt{OUTPUT: False} \\
            }
        }
    }
     \\ \\

    \theutterance \stepcounter{utterance}  
    & \multicolumn{4}{p{0.6\linewidth}}{
        \cellcolor[rgb]{0.561,0.737,0.561}{
            \makecell[{{p{\linewidth}}}]{
                \texttt{\tiny{[P1$\rangle$GM]}}
                \texttt{\{} \\
\texttt{  "input": {[}1, {-}1{]},} \\
\texttt{  "predicted\_output": 0,} \\
\texttt{  "current\_mode": "probing",} \\
\texttt{  "input\_rationale": "Since f(1,2)=3, f(2,3)=5, f(0,0)=0, f(1,1)=2, f(10,20)=30 are all True for addition, but f({-}1,1)=0 is False, the function is NOT simple addition. It behaves like addition for non{-}negative integers, but fails for negative ones. Let's test f(1, {-}1). If it's not 0, it might be something like bitwise OR, or max, or abs(a)+abs(b), or a*b, etc. Let's try a case that distinguishes these. If it were max(1, {-}1), output is 1. If it were abs(1)+abs({-}1), output is 2. If it were 1*{-}1, output is {-}1. If it were 1|{-}1, output is {-}1 (in two's complement) or depends on implementation. Let's try f(1, {-}1) with predicted output 1 (max).",} \\
\texttt{  "hypothesis": "unknown",} \\
\texttt{  "confidence\_score": 40} \\
\texttt{\}} \\
            }
        }
    }
    & & \\ \\

    \theutterance \stepcounter{utterance}  
    & & & \multicolumn{4}{p{0.6\linewidth}}{
        \cellcolor[rgb]{0.9,0.9,0.9}{
            \makecell[{{p{\linewidth}}}]{
                \texttt{\tiny{[P1$\langle$GM]}}
                \texttt{OUTPUT: False} \\
            }
        }
    }
     \\ \\

    \theutterance \stepcounter{utterance}  
    & \multicolumn{4}{p{0.6\linewidth}}{
        \cellcolor[rgb]{0.561,0.737,0.561}{
            \makecell[{{p{\linewidth}}}]{
                \texttt{\tiny{[P1$\rangle$GM]}}
                \texttt{\{} \\
\texttt{  "input": {[}1, {-}1{]},} \\
\texttt{  "predicted\_output": 2,} \\
\texttt{  "current\_mode": "probing",} \\
\texttt{  "input\_rationale": "Testing if the function is sum of absolute values: abs(1) + abs({-}1) = 2. This would explain why f(1,2)=3, f(2,3)=5, etc. (all positive) but f({-}1,1)=0 is false (since abs({-}1)+abs(1)=2 != 0).",} \\
\texttt{  "hypothesis": "unknown",} \\
\texttt{  "confidence\_score": 50} \\
\texttt{\}} \\
            }
        }
    }
    & & \\ \\

    \theutterance \stepcounter{utterance}  
    & & & \multicolumn{4}{p{0.6\linewidth}}{
        \cellcolor[rgb]{0.9,0.9,0.9}{
            \makecell[{{p{\linewidth}}}]{
                \texttt{\tiny{[P1$\langle$GM]}}
                \texttt{OUTPUT: True} \\
            }
        }
    }
     \\ \\

    \theutterance \stepcounter{utterance}  
    & \multicolumn{4}{p{0.6\linewidth}}{
        \cellcolor[rgb]{0.561,0.737,0.561}{
            \makecell[{{p{\linewidth}}}]{
                \texttt{\tiny{[P1$\rangle$GM]}}
                \texttt{\{} \\
\texttt{  "input": {[}{-}1, {-}1{]},} \\
\texttt{  "predicted\_output": 2,} \\
\texttt{  "current\_mode": "probing",} \\
\texttt{  "input\_rationale": "Testing the sum of absolute values hypothesis with two negative numbers. If f({-}1, {-}1) = 2, it confirms abs(a) + abs(b). Previous tests: f(1,2)=3 (1+2), f(2,3)=5 (2+3), f(0,0)=0 (0+0), f(1,1)=2 (1+1), f(10,20)=30 (10+20), f(1,{-}1)=2 (1+1). All consistent with abs(a)+abs(b). f({-}1,1)=0 was false because abs({-}1)+abs(1)=2 != 0.",} \\
\texttt{  "hypothesis": "def solution(a, b):\textbackslash{}n    \textbackslash{}"\textbackslash{}"\textbackslash{}"Returns the sum of the absolute values of a and b.\textbackslash{}"\textbackslash{}"\textbackslash{}"\textbackslash{}n    return abs(a) + abs(b)",} \\
\texttt{  "confidence\_score": 90} \\
\texttt{\}} \\
            }
        }
    }
    & & \\ \\

    \theutterance \stepcounter{utterance}  
    & & & \multicolumn{4}{p{0.6\linewidth}}{
        \cellcolor[rgb]{0.9,0.9,0.9}{
            \makecell[{{p{\linewidth}}}]{
                \texttt{\tiny{[P1$\langle$GM]}}
                \texttt{OUTPUT: True} \\
            }
        }
    }
     \\ \\

    \theutterance \stepcounter{utterance}  
    & \multicolumn{4}{p{0.6\linewidth}}{
        \cellcolor[rgb]{0.561,0.737,0.561}{
            \makecell[{{p{\linewidth}}}]{
                \texttt{\tiny{[P1$\rangle$GM]}}
                \texttt{\{} \\
\texttt{  "input": {[}{-}5, 10{]},} \\
\texttt{  "predicted\_output": 15,} \\
\texttt{  "current\_mode": "confirming",} \\
\texttt{  "input\_rationale": "Testing the hypothesis with mixed signs and larger magnitudes. abs({-}5) + abs(10) = 5 + 10 = 15. This further confirms the sum of absolute values hypothesis.",} \\
\texttt{  "hypothesis": "def solution(a, b):\textbackslash{}n    \textbackslash{}"\textbackslash{}"\textbackslash{}"Returns the sum of the absolute values of a and b.\textbackslash{}"\textbackslash{}"\textbackslash{}"\textbackslash{}n    return abs(a) + abs(b)",} \\
\texttt{  "confidence\_score": 95} \\
\texttt{\}} \\
            }
        }
    }
    & & \\ \\

    \theutterance \stepcounter{utterance}  
    & & & \multicolumn{4}{p{0.6\linewidth}}{
        \cellcolor[rgb]{0.9,0.9,0.9}{
            \makecell[{{p{\linewidth}}}]{
                \texttt{\tiny{[P1$\langle$GM]}}
                \texttt{OUTPUT: True} \\
            }
        }
    }
     \\ \\

    \theutterance \stepcounter{utterance}  
    & \multicolumn{4}{p{0.6\linewidth}}{
        \cellcolor[rgb]{0.561,0.737,0.561}{
            \makecell[{{p{\linewidth}}}]{
                \texttt{\tiny{[P1$\rangle$GM]}}
                \texttt{\{} \\
\texttt{  "input": {[}{-}100, {-}200{]},} \\
\texttt{  "predicted\_output": 300,} \\
\texttt{  "current\_mode": "confirming",} \\
\texttt{  "input\_rationale": "Final confirmation with large negative numbers. abs({-}100) + abs({-}200) = 100 + 200 = 300. This strongly supports the hypothesis that the function returns the sum of absolute values.",} \\
\texttt{  "hypothesis": "def solution(a, b):\textbackslash{}n    \textbackslash{}"\textbackslash{}"\textbackslash{}"Returns the sum of the absolute values of a and b.\textbackslash{}"\textbackslash{}"\textbackslash{}"\textbackslash{}n    return abs(a) + abs(b)",} \\
\texttt{  "confidence\_score": 99} \\
\texttt{\}} \\
            }
        }
    }
    & & \\ \\

    \theutterance \stepcounter{utterance}  
    & & & \multicolumn{2}{p{0.3\linewidth}}{
        \cellcolor[rgb]{0.9,0.9,0.9}{
            \makecell[{{p{\linewidth}}}]{
                \texttt{\tiny{[GM$|$GM]}}
                \texttt{Game Verdict: LOSS (turn limit reached)} \\
            }
        }
    }
    & & \\ \\

\end{supertabular}
}

\end{document}
